\documentclass[11pt,a4paper]{article}

\usepackage[utf8]{inputenc}
\usepackage[T1]{fontenc}
\usepackage[ngerman]{babel}
\usepackage[margin=2.5cm]{geometry}
\usepackage{amsmath}
\usepackage{booktabs}
\usepackage{xcolor}
\usepackage{listings}
\usepackage{hyperref}
\usepackage{enumitem}

\definecolor{codebg}{rgb}{0.95,0.95,0.95}
\lstset{
    basicstyle=\ttfamily\small,
    backgroundcolor=\color{codebg},
    frame=single,
    breaklines=true,
    columns=fullflexible,
    xleftmargin=0.3cm,
    xrightmargin=0.3cm,
    aboveskip=0.3cm,
    belowskip=0.3cm
}

\title{Beste Teillösungen für ungelöste Aufgaben\\ (DeepCoder)}
\author{Bemerkung zur Bachelorarbeit -- Steve Leonel Yomi Mbiakop}
\date{\today}

\sloppy

\begin{document}
\maketitle

\section{Ausgangslage}

Wie in der separaten Bemerkung \emph{``Warum liefern DeepCoder und
DreamCoder für ungelöste Aufgaben kein Kandidatenprogramm?''} gezeigt,
speicherte die DeepCoder-Suche (\texttt{deepcoder/search.py}) für eine
Aufgabe nur dann ein Programm, wenn dieses \textbf{alle} gegebenen
Beispiele exakt erfüllte. Scheiterte jeder während der Suche geprüfte
Kandidat an mindestens einem Beispiel, blieb \texttt{solution = None} --
und alle Strukturmetriken (OperationScore, PositionScore, OrderScore,
EditScore, PartialCorrectness) waren für diese Aufgabe rechnerisch immer
$0$, unabhängig davon, wie nah die Suche tatsächlich gekommen war.

Dieses Dokument beschreibt, wie DeepCoder so erweitert wurde, dass der
\textbf{beste gefundene Teiltreffer} zusätzlich gespeichert und für die
Metrikberechnung genutzt wird -- ausschließlich für DeepCoder (nicht für
DreamCoder, siehe Abschnitt \ref{sec:scope}).

\section{Lösung}

\subsection{\texttt{deepcoder/search.py}: Teiltreffer während der Suche verfolgen}

Neue Hilfsfunktion, die zählt, wie viele Beispiele ein Kandidat erfüllt
(statt wie \texttt{is\_solution} beim ersten Fehlschlag abzubrechen):

\begin{lstlisting}[language=Python]
def count_matches(program, examples):
    return sum(1 for inputs, output in examples
               if program(*inputs) == output)
\end{lstlisting}

In \texttt{dfs()} wird für jeden Kandidaten, der \emph{nicht} alle
Beispiele erfüllt, dessen Trefferzahl bestimmt. Übertrifft sie den
bisher besten Wert, wird der Kandidat gemerkt:

\begin{lstlisting}[language=Python]
ns = { 'nb_steps': 0,
       'solution': None,
       'gas' : gas,
       'best_partial': None,
       'best_partial_count': 0 }
...
try:
    if is_solution(p_base, examples):
        ns['solution'] = p_base
        return True
    match_count = count_matches(p_base, examples)
    if match_count > ns['best_partial_count']:
        ns['best_partial'] = p_base
        ns['best_partial_count'] = match_count
except (NullInputError, OutputOutOfRangeError):
    return
...
return ns['solution'], ns['nb_steps'], ns['best_partial'], ns['best_partial_count']
\end{lstlisting}

Findet sich nie ein Kandidat, der auch nur ein Beispiel trifft, bleibt
\texttt{best\_partial} bei \texttt{None} -- die Aufgabe fällt dann exakt
auf das alte Verhalten zurück (siehe Abschnitt \ref{sec:edgecase}). Der
alternative Suchmodus \texttt{sort\_and\_add()} wurde analog erweitert
und vergleicht den besten Teiltreffer über alle geprüften Kontexte hinweg.

\subsection{\texttt{scripts/solve-problems.py}: neue Ausgabespalten}

Für jede Aufgabe werden jetzt zwei zusätzliche Spalten geschrieben,
befüllt nur wenn die Aufgabe \emph{nicht} gelöst wurde:

\begin{itemize}[nosep]
    \item \texttt{best\_partial\_solution} -- Programmstring des besten
    Teiltreffers (oder leer).
    \item \texttt{best\_partial\_match\_fraction} -- Anteil erfüllter
    Beispiele (z.\,B. $0.4$ bei 2 von 5 Beispielen).
\end{itemize}

\subsection{\texttt{scripts/export\_results\_to\_csv.py}: Spalten durchreichen}

Reicht die beiden neuen Spalten unverändert von der \texttt{.h5}- in die
\texttt{.csv}-Datei durch (rückwärtskompatibel: ältere \texttt{.h5}-Dateien
ohne diese Spalten funktionieren weiterhin).

\subsection{\texttt{scripts/compute\_metrics.py}: Best-Effort-Metriken}

Neue Funktion, die entscheidet, welches Programm für den strukturellen
Vergleich herangezogen wird:

\begin{lstlisting}[language=Python]
def best_effort_program(row, has_partial_column):
    if row["solved"]:
        return row["solution"]
    if has_partial_column and is_present(row.get("best_partial_solution")):
        return row["best_partial_solution"]
    return None
\end{lstlisting}

Auf Basis dieses Programms werden fünf neue Spalten berechnet --
\texttt{best\_effort\_operation\_score}, \texttt{\_position\_score},
\texttt{\_order\_score}, \texttt{\_edit\_score} und
\texttt{\_partial\_correctness} -- mit denselben Metrik-Funktionen, die
bereits für gelöste Aufgaben verwendet werden. Die bisherigen Spalten
(\texttt{accuracy}, \texttt{partial\_correctness} usw.) bleiben
unverändert erhalten; die neuen Spalten kommen zusätzlich hinzu.

\subsection{Notebook: zwei neue Diagramme}

\texttt{deepcoder\_metric\_analysis.ipynb} zeigt jetzt zusätzlich zum
bestehenden Gesamtvergleich zwei eigene Balkendiagramme:
\begin{itemize}[nosep]
    \item \textbf{Nur gelöste Aufgaben} -- Strukturmetriken, ohne vs.
    mit Training.
    \item \textbf{Nur ungelöste Aufgaben} -- dieselben Metriken, aber auf
    Basis von \texttt{best\_effort\_*}, statt pauschal $0$.
\end{itemize}

\section{Sonderfall: gar kein Kandidat gefunden}
\label{sec:edgecase}

Manche Aufgaben bleiben auch mit dieser Erweiterung bei $0$: wenn die
Suche innerhalb des Budgets (\texttt{gas}) \emph{nie} ein Programm
gefunden hat, das auch nur ein einziges Beispiel erfüllt, bleibt
\texttt{best\_partial\_solution} leer, \texttt{best\_effort\_program}
liefert \texttt{None}, und alle \texttt{best\_effort\_*}-Metriken werden
korrekt zu $0$ (weil \texttt{tokenize\_program(None)} eine leere
Token-Liste ergibt -- dieselbe Rechnung wie zuvor bei einer leeren
\texttt{solution}). Kein Fehler, kein Absturz, keine künstliche
Verzerrung -- nur eine ehrliche Rückkehr zum alten Zustand für genau
diese Aufgaben.

\section{Validierung}

Testlauf auf \texttt{T=2\_test} (100 Aufgaben, gas\,=\,1000, ohne
Training):

\begin{center}
\begin{tabular}{lrr}
\toprule
& \textbf{vorher} & \textbf{nachher} \\
\midrule
Gelöste Aufgaben & 53/100 & 53/100 (unverändert) \\
Ungelöste Aufgaben mit Teiltreffer & -- & 25/47 \\
Ungelöste Aufgaben ganz ohne Kandidat & -- & 22/47 \\
Ø \texttt{partial\_correctness} (ungelöst) & 0{,}0000 & -- \\
Ø \texttt{best\_effort\_partial\_correctness} (ungelöst) & -- & 0{,}2556 \\
\bottomrule
\end{tabular}
\end{center}

Die Lösungsrate (\texttt{solved}) blieb exakt gleich -- die Erweiterung
verändert nicht, \emph{ob} eine Aufgabe als gelöst gilt, sondern liefert
für die verbleibenden Fälle zusätzliche, ehrliche Information über die
strukturelle Nähe der Suche zum Referenzprogramm.

\textbf{Konkretes Beispiel} (Aufgabe \texttt{LIST|COUNT,<0,0|TAKE,1,0},
ungelöst): der beste gefundene Teiltreffer erfüllt $2$ von $5$ Beispielen
($40\,\%$) und erhält dadurch einen echten \texttt{best\_effort\_*}-Score
statt einer pauschalen Null.

\section{Umfang: nur DeepCoder}
\label{sec:scope}

Diese Erweiterung wurde bewusst nur für DeepCoder umgesetzt. Dieser
Abschnitt begründet das nicht nur, sondern zeigt auch konkret, was für
DreamCoder nötig wäre, und warum die Entscheidung dagegen eine bewusste
Kosten-Nutzen-Abwägung war -- keine Verlegenheitslösung.

\subsection{Warum DreamCoder strukturell anders liegt}

Bei DeepCoder reichte eine reine Python-Änderung, weil die komplette
Suche in Python läuft (\texttt{deepcoder/search.py}). Bei DreamCoder
liegt die eigentliche Enumeration dagegen im kompilierten OCaml-Solver.
Wie in der separaten Bemerkung gezeigt, entscheidet dort
\texttt{solvers/task.ml} pro Kandidat binär:

\begin{lstlisting}[language=Caml]
let logLikelihood = tasks.(j).log_likelihood p in
if is_valid logLikelihood then begin
  Heap.add hits.(j) {hit_program = ...; hit_likelihood = logLikelihood; ...}
end
\end{lstlisting}

Ein Kandidat mit \texttt{logLikelihood = neg\_infinity} (scheitert an
mindestens einem Beispiel) besteht \texttt{is\_valid} nicht und wird
\emph{nie} in die Heap aufgenommen, die später als JSON an Python
zurückgegeben wird (\texttt{solvers/solver.ml},
\texttt{export\_frontiers}). Der Python-Wrapper
(\texttt{solveForTask\_ocaml} in \texttt{dreamcoder/enumeration.py})
bekommt also gar keine Kandidaten mit Teiltreffern zu Gesicht -- sie
existieren zu keinem Zeitpunkt außerhalb des OCaml-Prozesses.

\subsection{Was die Umsetzung technisch erfordern würde}

Um das DeepCoder-Prinzip (bester Teiltreffer statt Verwerfen) auch für
DreamCoder umzusetzen, wären mindestens vier zusammenhängende Änderungen
nötig:

\begin{enumerate}[nosep]
    \item \textbf{\texttt{solvers/task.ml}, \texttt{supervised\_task}:}
    Eine zweite, graduelle Bewertungsfunktion ergänzen, die zählt, wie
    viele Beispiele \texttt{p} erfüllt (analog zu \texttt{count\_matches}
    in DeepCoder), statt beim ersten Fehlschlag mit \texttt{log 0.0}
    abzubrechen.
    \item \textbf{\texttt{solvers/task.ml}, \texttt{enumerate\_for\_tasks}:}
    Pro Aufgabe einen zweiten, von \texttt{hits} getrennten Merker für
    den bisher besten Teiltreffer einführen (die \texttt{is\_valid}-Prüfung
    darf für \texttt{hits} unverändert bleiben, damit sich an gelösten
    Aufgaben nichts ändert).
    \item \textbf{\texttt{solvers/solver.ml}, \texttt{export\_frontiers}:}
    Das JSON-Antwortformat um ein optionales Feld für den besten
    Teiltreffer je Aufgabe erweitern.
    \item \textbf{\texttt{dreamcoder/enumeration.py},
    \texttt{solveForTask\_ocaml}:} Das neue Feld einlesen und (ähnlich wie
    bei DeepCoder) getrennt von der regulären \texttt{Frontier} als
    \texttt{best\_partial}-Information an die Auswertungs-Pipeline
    weiterreichen.
\end{enumerate}

Jede dieser vier Änderungen ist für sich genommen nicht kompliziert --
das Muster ist ja bei DeepCoder bereits erfolgreich erprobt. Der
eigentliche Aufwand liegt nicht in der Logik, sondern in der
\textbf{Bauumgebung}, siehe unten.

\subsection{Konkret geprüfter Befund: keine Build-Umgebung vorhanden}

\begin{lstlisting}[language=bash]
$ which ocaml ocamlfind dune opam
(nichts gefunden)
$ opam --version
bash: opam: command not found
$ find / -iname '*opam*'
/etc/apparmor.d/opam        (AppArmor-Profil, kein Programm)
/usr/share/vim/.../opam.vim (Vim-Syntaxdatei, kein Programm)
\end{lstlisting}

In der WSL-Umgebung ist aktuell \emph{kein} OCaml-Build-Toolchain
installiert. Die vorhandenen Binaries \texttt{solver} und
\texttt{compression} sind vorgefertigt (Änderungsdatum 10.\,Juli, aus der
ursprünglichen Repository-Einrichtung) und lokal nicht neu baubar. Das
Build-System selbst (\texttt{solvers/Makefile}) verrät zusätzlich, mit
welchem -- inzwischen unüblichen -- Werkzeug gebaut wurde:

\begin{lstlisting}[language=bash]
c = corebuild -pkg yojson -pkg re2 -tag unsafe -tag "optimize(3)"
\end{lstlisting}

\texttt{corebuild} ist ein älteres, auf Jane Streets \texttt{Core}-
Bibliothek zugeschnittenes Build-Werkzeug (Vorgänger des heute üblichen
\texttt{dune}). Zusätzliche Abhängigkeit: \texttt{re2}, eine
OCaml-Anbindung an Googles gleichnamige C++-Regex-Bibliothek -- solche
nativen Bindings sind erfahrungsgemäß besonders versionsempfindlich beim
Neuaufsetzen. Es existiert weder ein Dockerfile noch ein Setup-Skript im
Repository, das den ursprünglichen Build-Prozess dokumentiert oder
reproduzierbar macht.

\subsection{Kosten-Nutzen-Abwägung}

Vor der eigentlichen Code-Änderung (Abschnitt 5.2) stünde also erst der
Aufbau einer kompletten OCaml-Toolchain: \texttt{opam} installieren,
einen zur \texttt{Core}-Version passenden OCaml-Switch anlegen,
\texttt{core}, \texttt{yojson} und \texttt{re2} samt alle transitiven
Abhängigkeiten installieren, und danach verifizieren, dass sich der
\emph{unveränderte} Code überhaupt noch baut, bevor die eigentliche
Änderung überhaupt begonnen werden kann. Realistischer Aufwand: 1\,+\,
Stunden, mit ungewissem Ausgang -- gerade \texttt{re2} ist ein bekannter
Quell für Build-Fehler bei neu aufgesetzten OCaml-Umgebungen.

Angesichts dieses Aufwands-Nutzen-Verhältnisses -- und angesichts der
Erfahrung aus der heutigen Arbeit, wie instabil die WSL2-Umgebung bereits
bei reinen Python-Läufen war (Speicherabstürze, siehe separate
Bemerkung) -- wurde bewusst entschieden, diese Erweiterung
\textbf{nicht} umzusetzen. Die vier oben skizzierten Schritte bleiben
als konkret ausformulierter, sofort umsetzbarer nächster Schritt
festgehalten, falls zu einem späteren Zeitpunkt eine funktionierende
OCaml-Umgebung zur Verfügung steht.

\textbf{Wichtig für die Einordnung der Ergebnisse:} Diese Entscheidung
schränkt \emph{nicht} die Gültigkeit der bisherigen DreamCoder-Ergebnisse
ein. Sie bedeutet lediglich, dass für DreamCoders ungelöste Aufgaben --
anders als bei DeepCoder -- weiterhin nur der binäre Zustand
``gelöst/nicht gelöst'' und keine graduelle Teilkorrektheit verfügbar
ist.

\end{document}
